\section{Motivation}

Computer science students meet abstraction every day: an \emph{interface} specifies which operations a
data type must support and which laws they must obey, and any code written against the interface
automatically works for every implementation. This chapter does exactly that to arithmetic. The arithmetic
of the complex numbers we have just met — and of $\Q$ and $\R$ before them — obeys the same short list of
rules: operations that are associative, have a neutral element, and admit inverses. Abstracting this list
leads to two of the most basic algebraic ``interfaces'': \emph{groups} and \emph{fields}. Every theorem we
prove from the axioms alone is proved once, for all implementations simultaneously: for real numbers, for
complex numbers, for residues modulo a prime, for invertible matrices.

This is not abstraction for its own sake — the ``exotic'' implementations are precisely the ones running
your digital life:

\begin{itemize}
\item \textbf{Cryptography.} Whenever your browser shows a padlock, a key exchange has just been performed
in a large finite group (Diffie--Hellman in $\Z_p\setminus\{0\}$, or on an elliptic curve). RSA encryption
is arithmetic modulo $n$, and its security rests on how hard certain computations in these groups are. The
finite fields $\Z_p$ defined in this chapter are the entry ticket to all of modern cryptography.
\item \textbf{Error-correcting codes.} QR codes, scratched CDs, hard-disk RAID arrays and photographs sent
back by space probes all survive damage thanks to Reed--Solomon codes — linear algebra over a finite field
with $256$ elements. The AES cipher that encrypts your Wi-Fi traffic computes in the same field.
\item \textbf{Bits form a field.} The smallest field, $\Z_2=\{0,1\}$, is the arithmetic of a single bit:
its addition is the XOR gate and its multiplication is the AND gate. Linear algebra over $\Z_2$ is a
standard tool in hashing, checksums and cryptanalysis.
\end{itemize}

For this course, though, the most important reason is different: fields are exactly the sets of
\emph{scalars} over which vector spaces — the central objects of linear algebra — will be built in the
following chapters. Whenever a definition below feels dry, remember: we are writing the interface, and the
rest of this book is the code that runs against it.

\section{Groups}

\begin{definition}
A \emph{group} is a set $G$ together with a binary operation $*\colon G\times G\to G$ satisfying:
\begin{itemize}
\item \textbf{(Associativity)} $(a*b)*c = a*(b*c)$ for all $a,b,c\in G$;
\item \textbf{(Identity)} there is an element $e\in G$ such that $a*e=e*a=a$ for all $a\in G$;
\item \textbf{(Inverses)} for every $a\in G$ there is an element $a^{-1}\in G$ such that
$a*a^{-1}=a^{-1}*a=e$.
\end{itemize}
A group is called \emph{abelian} (or \emph{commutative}) if, in addition, $a*b=b*a$ for all $a,b\in G$.
\end{definition}

\begin{remark}
When the operation is written additively (denoted $+$) the identity is written $0$ and the inverse of $a$ is
written $-a$. When it is written multiplicatively (denoted $\cdot$) the identity is written $1$ and the
inverse of $a$ is written $a^{-1}$ or $\tfrac{1}{a}$.
\end{remark}

Before listing examples we introduce one construction that will recur throughout the course:
\emph{modular arithmetic}.

\begin{definition}[Modular arithmetic]
Fix a positive integer $n$. For an integer $a$ we write $a \bmod n$ for the remainder of the division of
$a$ by $n$: the unique number $r\in\{0,1,\ldots,n-1\}$ such that $a = qn + r$ for some $q\in\Z$. Two
integers $a$ and $b$ are \emph{congruent modulo $n$} if they leave the same remainder, i.e.\ if $n$
divides $a-b$. On the set $\Z_n = \{0,1,\ldots,n-1\}$ we define \emph{addition and multiplication modulo
$n$}:
\[
a +_n b = (a+b) \bmod n, \qquad a \cdot_n b = (a\cdot b) \bmod n.
\]
\end{definition}

\begin{example}
$17 \bmod 4 = 1$ (since $17 = 4\cdot 4+1$), $\ 4 \bmod 17 = 4$, and $(-2)\bmod 5 = 3$ (since
$-2 = (-1)\cdot 5 + 3$; the remainder is always taken in $\{0,\ldots,n-1\}$). In $\Z_6$ we have
$4 +_6 5 = 3$ and $4 \cdot_6 5 = 2$.
\end{example}

\begin{example}
\hfill
\begin{itemize}
\item $(\Z,+)$, $(\Q,+)$, $(\R,+)$ and $(\C,+)$ are abelian groups; the identity is $0$ and the inverse of
$a$ is $-a$.
\item $(\Q\setminus\{0\},\cdot)$, $(\R\setminus\{0\},\cdot)$ and $(\C\setminus\{0\},\cdot)$ are abelian
groups; the identity is $1$ and the inverse of $a$ is $\tfrac{1}{a}$.
\item $(\Z,\cdot)$ is \emph{not} a group: apart from $\pm 1$, no integer has a multiplicative inverse in
$\Z$.
\item For a fixed $n$, the set $\Z_n=\{0,1,\ldots,n-1\}$ with addition modulo $n$ is an abelian group.
\end{itemize}
\end{example}

\begin{example}[A finite multiplicative group]
For a prime $p$, the non-zero residues $\Z_p\setminus\{0\}=\{1,2,\ldots,p-1\}$ form an abelian group under
multiplication modulo $p$. For $p=5$, for example, $2\cdot 3 = 6 \equiv 1 \pmod 5$, so $3$ is the inverse of
$2$. Inverses allow us to solve linear equations exactly as in $\Q$ or $\R$: in $\Z_7$ the equation
$2\cdot_7 x = 1$ has the unique solution $x = 4$ (indeed $2\cdot 4 = 8 \equiv 1 \pmod 7$), and
$2 \cdot_7 x = 5$ is solved by multiplying both sides by $4$: $x = 4\cdot_7 5 = 6$.
\end{example}

\begin{example}[A non-abelian group]
The set of all invertible $2\times 2$ real matrices with matrix multiplication forms a group, called the
\emph{general linear group} $GL_2(\R)$ (matrices and their multiplication are treated in detail in
Chapter~\ref{ch:matrices}). It is \emph{not} abelian; for instance
\[
\begin{pmatrix} 1 & 1\\ 0 & 1\end{pmatrix}
\begin{pmatrix} 1 & 0\\ 1 & 1\end{pmatrix}
=\begin{pmatrix} 2 & 1\\ 1 & 1\end{pmatrix}
\neq
\begin{pmatrix} 1 & 1\\ 1 & 2\end{pmatrix}
=\begin{pmatrix} 1 & 0\\ 1 & 1\end{pmatrix}
\begin{pmatrix} 1 & 1\\ 0 & 1\end{pmatrix}.
\]
\end{example}

\begin{example}[The power-set group]
Let $X$ be a set and let $2^X$ denote the set of all its subsets. The \emph{symmetric difference}
\[
A \bigtriangleup B = (A\cup B)\setminus (A\cap B)
\]
(the set of elements belonging to exactly one of $A$ and $B$) makes $(2^X,\bigtriangleup)$ an abelian
group: the identity element is $\emptyset$, and every element is its own inverse, since
$A \bigtriangleup A = \emptyset$. Identifying a subset $A\subseteq \{1,\ldots,n\}$ with its characteristic
bit vector, $\bigtriangleup$ becomes the bitwise XOR of bit strings.
\end{example}

\begin{fact}
In any group $G$ the identity element is unique, and each element has a unique inverse.
\end{fact}

\begin{proof}
If $e$ and $e'$ are both identities, then $e=e*e'=e'$. If $b$ and $b'$ are both inverses of $a$, then
\[
b = b*e = b*(a*b') = (b*a)*b' = e*b' = b'. \qedhere
\]
\end{proof}

\begin{fact}[Cancellation laws]
In any group $G$: if $a*b = a*c$, then $b=c$; similarly, if $b*a = c*a$, then $b=c$.
\end{fact}

\begin{proof}
Multiplying $a*b=a*c$ by $a^{-1}$ on the left gives $a^{-1}*(a*b) = a^{-1}*(a*c)$; by associativity
$(a^{-1}*a)*b = (a^{-1}*a)*c$, that is $e*b = e*c$, so $b=c$. The second law is proved symmetrically.
\end{proof}

\section{Fields}

A field is a set carrying \emph{two} compatible operations, addition and multiplication, each of which behaves
like an abelian group (with multiplication restricted to the non-zero elements), tied together by
distributivity. The fields $\Q$ and $\R$ are the familiar prototypes.

\begin{definition}
A set $\K$ together with two binary operations $+$ and $\cdot$ is called a \emph{field} if it satisfies the
following conditions.
\begin{itemize}
\item for any $a,b,c\in \K$, $\ a+b=b+a\ $ and $\ a+(b+c)=(a+b)+c$;
\item there is an element $0\in \K$ such that $a+0=0+a=a$ for any $a\in \K$;
\item for any $a\in \K$ there exists $-a\in \K$ such that $a+(-a)=(-a)+a=0$;
\item for any $a,b,c\in \K$, $\ a\cdot b = b\cdot a\ $ and $\ a\cdot(b\cdot c)=(a\cdot b)\cdot c$;
\item there is an element $1\in \K$ such that $a\cdot 1=1\cdot a=a$ for any $a\in \K$;
\item for any $a\in \K$ with $a\neq 0$ there is $\tfrac{1}{a}\in \K$ such that
$a \cdot \tfrac{1}{a} = \tfrac{1}{a}\cdot a = 1$;
\item for any $a,b,c\in \K$, $\ (a+b)\cdot c = a\cdot c + b\cdot c$.
\end{itemize}
\end{definition}

\begin{remark}
Equivalently: $(\K,+)$ is an abelian group with identity $0$, the non-zero elements
$(\K\setminus\{0\},\cdot)$ form an abelian group with identity $1$, and multiplication distributes over
addition.
\end{remark}

\begin{example}
$\Q$, $\R$ and $\C$, each with their usual addition and multiplication, are fields.
\end{example}

\begin{example}
Consider the set $\Z_2=\{0,1\}$ with addition and multiplication given by
\[
\begin{array}{c|cc} + & 0 & 1 \\ \hline 0 & 0 & 1 \\ 1 & 1 & 0 \end{array}
\qquad\qquad
\begin{array}{c|cc} \cdot & 0 & 1 \\ \hline 0 & 0 & 0 \\ 1 & 0 & 1 \end{array}
\]
Then $\Z_2$ with these operations is a field (the smallest one).
\end{example}

\begin{example}
More generally, consider $\Z_n$ with addition and multiplication modulo $n$.
\begin{itemize}
\item $\Z_4$ is \emph{not} a field: the products $2\cdot 0 = 0$, $2\cdot 1 = 2$, $2\cdot 2 = 0$ and
$2\cdot 3 = 2$ (all modulo $4$) never equal $1$, so the element $2$ has no multiplicative inverse.
\item $\Z_5$ \emph{is} a field: modulo $5$ we have $1\cdot 1 = 1$, $2\cdot 3 = 1$ and $4\cdot 4 = 1$, so
every non-zero element has an inverse.
\end{itemize}
\end{example}

\begin{theorem}
$\Z_n$, with addition and multiplication modulo $n$, is a field if and only if $n$ is a prime number.
\end{theorem}

\begin{example}[A field between $\Q$ and $\R$]
The set
\[
\Q(\sqrt 2) = \{a+b\sqrt 2\mid a,b\in \Q\}\subseteq \R
\]
with the usual addition and multiplication is a field. The only non-obvious axiom is the existence of
multiplicative inverses, and it follows from the same trick we used for complex numbers: if
$a+b\sqrt2\neq 0$, then
\[
\frac{1}{a+b\sqrt 2} = \frac{a-b\sqrt2}{(a+b\sqrt2)(a-b\sqrt2)}
= \frac{a}{a^2-2b^2}-\frac{b}{a^2-2b^2}\,\sqrt 2,
\]
where $a^2-2b^2\neq 0$ because $\sqrt2$ is irrational.
\end{example}

\begin{warning}
Not every set with an addition and a multiplication is a field. For example, $\Z$ with the standard operations
is \emph{not} a field, because only $1$ and $-1$ have multiplicative inverses in $\Z$.
\end{warning}

\noindent In the chapters that follow, $\K$ always denotes a field; the reader may safely keep $\R$ or $\C$ in
mind throughout.
